% ECONOMICS_PROBLEM: {"id": "F22-536", "title": "Migration, remittances and structural transformation", "jel_code": "F22", "source_id": "OP-0038"}
% !TeX program = xelatex
\documentclass[11pt,a4paper]{article}
\usepackage[margin=23mm]{geometry}
\usepackage{fontspec}
\setmainfont{Latin Modern Roman}
\usepackage{amsmath,amssymb,mathtools}
\usepackage{xcolor,enumitem,booktabs,longtable,array,xurl,hyperref}
\definecolor{muted}{HTML}{53636C}
\hypersetup{colorlinks=true,urlcolor=blue,linkcolor=blue}
\setlength{\parindent}{0pt}
\setlength{\parskip}{5pt}
\newcommand{\R}{\mathbb R}
\newcommand{\E}{\mathbb E}
\newcommand{\N}{\mathbb N}
\newcommand{\ind}{\mathbf 1}
\newcommand{\Var}{\operatorname{Var}}
\newcommand{\Cov}{\operatorname{Cov}}
\newcommand{\supp}{\operatorname{supp}}
\DeclareMathOperator*{\argmin}{arg\,min}
\DeclareMathOperator*{\argmax}{arg\,max}
\newcommand{\entrymeta}[3]{{\sffamily\small\color{muted}\textbf{\texttt{#1}}\quad|\quad #2\par #3\par}}
\newcommand{\Problemref}[1]{\texttt{#1}}
\newcommand{\JELref}[2]{\texttt{#2} (JEL \texttt{#1})}
\begin{document}
\section*{Migration, remittances and structural transformation}
\textbf{Problem ID:} \texttt{F22-536}\quad \textbf{JEL:} \texttt{F22}\par
% BEGIN ECONOMICS STATEMENT
\label{problem:OP-0038}
\label{atlas:prob:D8}
\noindent\textbf{Original atlas ID:} D8 (entry 38). \textbf{Classification:} Editorial migration-policy and origin-transformation problem.

\noindent\textbf{Original atlas question:} Does migration accelerate development through income and knowledge or delay local productive transformation?

\subsection*{Definitions and assumptions}
Fix origin communities, baseline residents, possible destinations, and horizon $H$. Let $\pi$ specify migration eligibility, assistance, work rights, and coverage. Migration itself is endogenous unless directly assigned. Let $Y_{it}(\pi)$ be real earnings or consumption of baseline resident $i$, tracked abroad as well as at origin. Let $Q_{mt}(\pi)$ be origin-community value added, $S_{mt}(\pi)$ sectoral employment shares, and $R_{mt}(\pi)$ remittance receipts under a consistent definition. A spatial model $M$ includes worker selection, household bargaining, credit, knowledge transfer, exchange rates, firms, and local prices; specify its equilibrium and selection rule. Treatment may interfere within and between communities.
\subsection*{Formal research question}
For every $t\in\{0,\ldots,H\}$, define
\begin{align*}
 \Delta Y_t&=\mathbb E[Y_{it}(\pi)-Y_{it}(\pi_0)],\\
 \Delta V_t&=\mathbb E[V_{mt}(\pi)-V_{mt}(\pi_0)],
                  \qquad V\in\{Q,S,R\}.
\end{align*}
For supported policy coverage and baseline groups, identify or bound
\[
 \Theta_H(\pi,\pi_0)
    =(\Delta Y_t,\Delta Q_t,\Delta S_t,\Delta R_t)_{t=0}^H.
\]
Determine when gains to the baseline cohort coexist with durable growth in origin value added, rather than treating these as the same outcome. To separate remittance, labor-supply, and knowledge pathways, define component-specific interventions within $M$, account for their interactions, and identify which additional assignments or measurements discriminate their predicted effects.
\subsection*{Interpretation and scope}
Higher consumption at origin can arise from remittance transfers without an increase in local production. Migrants' earnings belong to the baseline-cohort outcome even though they are not origin GDP. Conversely, outcomes among remaining residents can change through selective departure. A migration lottery or encouragement identifies its offer effect; migration effects require compliance assumptions and an exclusion restriction. Exchange-rate shocks among existing migrant households chiefly identify a financial margin, and may also affect trade, destination labor conditions, or return decisions. They do not randomize who migrated. Small-scale encouragement can leave destination wages and origin prices approximately unchanged while large-scale policy changes them. Household splitting, currency conversion, and remittance costs must be measured consistently when comparing welfare across locations.
\subsection*{Source audit and status}
[S1] provides exchange-rate evidence, [S2] randomized seasonal-migration encouragement, and [S3] a review of large potential gains from emigration. These concern different margins and cannot be combined as if they estimate one treatment. The origin-transformation and equilibrium mechanism question is editorial. Existing local effects are acknowledged, and no globally unresolved sign or universal 2026 frontier is claimed.

\subsection*{References}
\begin{enumerate}[label=\textnormal{[S\arabic*]},leftmargin=*]
\item Dean Yang (2008). \emph{International Migration, Remittances and Household Investment: Evidence from Philippine Migrants' Exchange Rate Shocks}. The Economic Journal 118(528), 591--630. \url{https://doi.org/10.1111/j.1468-0297.2008.02134.x}. \textit{Locator:} Primary publisher, working-paper, or author record: bibliographic information and abstract; full-text theorem verification is not claimed. \textit{Role:} Exchange-rate variation among migrant households; identifies a remittance-related margin under exclusions.
\item Gharad Bryan, Shyamal Chowdhury, and Ahmed Mushfiq Mobarak (2014). \emph{Underinvestment in a Profitable Technology: The Case of Seasonal Migration in Bangladesh}. Econometrica 82(5). \url{https://doi.org/10.3982/ECTA10489}. \textit{Locator:} Primary publisher, working-paper, or author record: bibliographic information and abstract; full-text theorem verification is not claimed. \textit{Role:} Randomized encouragement of seasonal migration; local evidence on migration constraints.
\item Michael A. Clemens (2011). \emph{Economics and Emigration: Trillion-Dollar Bills on the Sidewalk?}. Journal of Economic Perspectives 25(3), 83--106. \url{https://doi.org/10.1257/jep.25.3.83}. \textit{Locator:} Primary publisher, working-paper, or author record: bibliographic information and abstract; full-text theorem verification is not claimed. \textit{Role:} Review of global gains from migration; not a direct origin-community scale-up estimate.
\end{enumerate}

\subsection*{Common conventions: Source mathematical conventions}

Unless an entry states otherwise, agent and firm sets are finite. State and action spaces are measurable, strategies and outcome maps are measurable, probability laws are specified, and expectations exist for the targets under discussion. Infinite-horizon discount factors lie in $(0,1)$ and discounted objectives are finite. Norms, tolerances and complexity input sizes must be specified for computational comparisons. Constants or primitives that appear locally are inputs, not empirical facts asserted by this catalogue.

\textbf{Identification.} A statistical model $m\in\mathcal M$ implies an observable law $P_m$ and target $\theta(m)$. At law $P$, its identified set is
\[
\Theta(P;\mathcal M)=\{\theta(m):m\in\mathcal M,\ P_m=P\}.
\]
Point identification holds when this set is a singleton, or equivalently when $P_{m_1}=P_{m_2}$ implies $\theta(m_1)=\theta(m_2)$ for all compatible models. Sharp bounds mean bounds attained, or approached when the boundary is not attained, by compatible models. Writing this set is a definition, not a solution: the substantive task is to establish informative restrictions and derive or estimate the set. An unrestricted-confounding model may give only trivial bounds.

\textbf{Inference.} Identification, estimability and valid uncertainty quantification are separate questions. Fix $\alpha\in(0,1)$. A confidence procedure $C_n$ over a declared law class $\mathcal P$ has asymptotic uniform coverage at level $1-\alpha$ when
\[
\liminf_{n\to\infty}\inf_{P\in\mathcal P}\Pr_P\{\theta(P)\in C_n\}\geq1-\alpha.
\]
For set-identified targets the coverage event must be defined separately, for example coverage of the full identified set. If an entry uses identification-indexed models instead of uniquely indexed laws, the same coverage convention is applied over those models. Pointwise limits do not establish uniform validity, and asymptotic validity does not guarantee finite-sample precision. Sample sizes, dependence structures and weak-identification sequences are part of each application.

\textbf{Counterfactuals.} $Y_i(a)$ is a potential outcome under a specified intervention, including its timing, implementation and versions. It is not $Y_i$ conditional on voluntarily choosing $a$. A full intervention vector permits interference; shorthand scalar treatment notation does not silently impose no interference. Assignment, selection, attrition, anticipation and measurement must be specified. A claim about an instrument's local effect is not automatically a population policy effect.

\textbf{Economic equilibrium.} A counterfactual model includes best responses, constraints, feasibility, market clearing, state transitions and expectations consistent with the stated information. When several equilibria exist, either a declared selector is an input or the target is set-valued. Comparisons across policy regimes must use the stated selection convention. Reduced-form relationships are not presumed invariant after an equilibrium-changing intervention.

\textbf{Optimization and welfare.} Optimization is over the stated feasible set. If attainment is not established, $\sup$ or $\inf$ and $\varepsilon$-optimal choices replace an asserted maximizer or minimizer. For example, with value $V=\sup_{a\in\mathcal A}W(a)<\infty$, an $\varepsilon$-optimal set is $\{a\in\mathcal A:W(a)\geq V-\varepsilon\}$ for $\varepsilon>0$. Benefits, costs, risk and health or environmental outcomes can be aggregated only using declared units and conversion weights. Interpersonal utility weights, the treatment of fiscal transfers and foreign profits, discounting, and the behavioral welfare criterion are explicit inputs. A comparative-static derivative additionally requires existence and appropriate differentiability of the selected counterfactual.

\textbf{Sources.} Local labels [S1], [S2], and so forth restart in every question. Locators identify the relevant abstract, model, section or result at the level actually checked; a bibliographic or abstract-level verification is not represented as inspection of an entire proof. Working papers are distinguished from later publications and changing versions. Source summaries are paraphrased. All URL checks and editorial formulations refer to the audit date stated above.
% END ECONOMICS STATEMENT
\end{document}
